\section{A common target through principal \texorpdfstring{\(W\)}{W}-algebras}\label{sec:walg}

The pairings of \eqref{eq:L3} take values in Whittaker coinvariants of
Kac--Moody modules.  We identify these targets for the two inverse shifted
levels.  At a single point this is categorical Feigin--Frenkel duality.
The assertion needed here is stronger: the identification must respect
moving points, collisions, and the crystal structure on the parameter
scheme.  We construct it on compact generators and prove the necessary
base-change assertion for their mapping complexes.

Throughout this section \(S\) is a smooth affine coordinate chart for a
finite tuple of points of the curve, \(R=\mathcal O(S)\), and \(u=P(z)\)
is the monic equation of their effective divisor in the coordinate
\(z\).  Write \(\mathcal O\) for its completed regular ring and
\(\mathcal K=\mathcal O[u^{-1}]\).  All residues are total residues with
respect to \(dz\).  We use the Kac--Moody category \(\mathcal M\), its
Whittaker coinvariants \(Q_S\), and the passage \(b\) to the underlying
\(\operatorname{QCoh}(S)\)-linear category from Section~\ref{sec:loop}.
A superscript \(0\) denotes this last passage, as in \(Q_S^0=bQ_S\).
The definitions are local on \(S\); the coordinate and crystal comparisons
below will make the resulting equivalence independent of this chart.

\subsection{Compact open approximations}

For a positive root \(\alpha\), let \(h_\alpha\) be its height.  The
adolescent Whittaker subgroup \(I_i\), for \(i\geq1\), has the following
root and Cartan orders:
\begin{equation}\label{eq:L4}
\begin{array}{c|ccc}
 &e_\alpha&\mathfrak t&f_\alpha\\ \hline
 I_i&-i h_\alpha&i&i(h_\alpha+1).
\end{array}
\end{equation}
Here an order \(m\) means coefficients in \(u^m\mathcal O\).
Write \(N_i\) for the positive-root part and \(B_i\) for the
Cartan and negative-root part.  These are pro-unipotent groups, with
product coordinates, and \(I_i=N_iB_i\).  The character \(\chi\)
extends to \(I_i\) by zero on \(B_i\).

The multiplicative twisting splits over these groups with the Lie
splitting specified by the standard cocycle.  Indeed the order bounds
make that cocycle zero.  At integral depth the splitting is the
conjugated arc splitting.  Its possible Cartan correction vanishes on
positive-order Cartan coefficients.  We will also use intermediate
rational orders.  To see the splitting there without assuming that
\(\rho\) is an integral cocharacter, give the root piece of signed height
\(l\) the order \(m+\lambda l\).  The group preserves the corresponding
rationally filtered arc lattice in the adjoint representation.  Its
order-zero action is upper unipotent in root height, and all other
pieces increase order.  The determinant of its action on every finite
lattice difference is therefore canonically trivial.  The differential
of this trivialization has zero root correction by weights and zero
Cartan correction because positive-order Cartan modes have zero trace
on the successive order quotients.  Passing through the finite
nilpotent quotients gives the asserted splitting.

\begin{proposition}[Adolescent presentation]\label{walg:adolescent}
There is a canonical presentation
\begin{equation}\label{eq:L5}
 Q_S\simeq\underset{i\geq1}{\operatorname{colim}}P_i,
 \qquad P_i=\mathcal M^{I_i,\chi},
\end{equation}
whose transition functors are star averaging with character.  They are
fully faithful and have continuous \(D(S)\)-linear right adjoints.
After passage to \(b\), they preserve compact objects.  These
assertions hold over collision loci and after the base changes of Section~\ref{sec:loop}.
The same presentation and transition assertions apply with
\(a_S\) in place of \(\mathcal M\).
\end{proposition}

\begin{proof}
First fix \(N_i\).  Its invariant category is generated by objects with
additional \(B_j\)-equivariance for sufficiently large \(j\geq i\).
In fact, the underlying category is generated by objects with deep
congruence equivariance.  For such an object choose \(j\) so that
\(B_j\) lies in that congruence subgroup.  Product coordinates identify
its \(N_i\)-averaging with its \(N_iB_j\)-averaging.  Forgetting the
additional equivariance is fully faithful and has continuous right
adjoint by Proposition~\ref{loop:relative-equivariance}.  Thus these
subcategories exhaust the \(N_i\)-invariant category.  Increasing
\(i\) commutes with these exhaustions after increasing \(j\), by the
same averaging formula.  The diagonal depths are cofinal.  This proves
\eqref{eq:L5}; bounded loop supports give the identical argument for
\(a_S\).

It remains to prove the transition assertion relatively.  We give the
Fourier argument underlying the adolescent construction
\cite[Section~2]{RaskinW}, retaining the divisor coefficients to account
for collisions.  Interpolate the orders using rational slopes
\(\lambda>1\).  At a break the new orders are
\[
 a_l=-\lfloor\lambda l\rfloor\quad(l>0),\qquad
 b_r=\lfloor\lambda(r+1)\rfloor\quad(r\geq0),
\]
where \(r=0\) denotes the Cartan part.  The old orders immediately
before this break are
\[
 a_l^0=a_l+\mathbf1_{\lambda l\in\mathbb Z},\qquad
 b_r^0=b_r-\mathbf1_{\lambda(r+1)\in\mathbb Z}.
\]
Let \(J\) use the new positive orders and the old nonpositive orders,
and let \(K\subset J\) use the old positive orders and the new
nonpositive orders.  Then \(K\) is normal, and its character is
\(J\)-invariant.  Here are the order checks.  In positive roots they
follow from the floor inequality
\(\lfloor x\rfloor+\lfloor y\rfloor\leq\lfloor x+y\rfloor\);
when the sum is an integral break, either both summands are integral
or their fractional parts add to one.  A sum of two old nonpositive
orders reaches the new order, since \(\lambda>1\).  For a mixed
bracket of positive height \(l\) and nonpositive height \(-r\),
\(a_l+b_r^0\) reaches the old nonpositive order if \(l\leq r\),
and can fail to reach the new order only when both participating
breaks are integral.  If \(l>r\), it already reaches the old positive
order of height \(l-r\).  A simple-root residue can occur only when
\(l-r=1\), both breaks are integral, and both coefficients are among
those changed at the break.  This verifies normality and invariance
of the character on \(K\).  In the Cartan case the orders also exclude
a residue in the central Lie cocycle.  Exponentiating the finite
nilpotent quotients gives the group assertions.

Modulo \(K\), denote the new/old positive difference by \(U\) and the
old/new nonpositive difference by \(V\).  The group \(U\) is unipotent,
and \(V\) is a normal vector group.  Relative Fourier transform for
\(V\), applied to the category of \((K,\chi)\)-equivariant objects,
identifies the old invariant category with the fiber over zero in
\(V^*\).  The new condition is \(U\)-equivariance.  Its action on
\(V^*\) is affine: exchanging the two factors in multiplication
introduces both conjugation on \(V\) and the residual character of
\(K\); Fourier transform turns the latter into translation.  Thus
this description also accounts for the character twisting of the
quotient action.

The orbit of zero is free and closed over \(S\).  To check this,
order coordinates by root height.  The linear term from a height
\(l\) coordinate of \(U\) to the height \(l-1\) component of
\(V^*\) is the pairing
\[
 (x,y)\longmapsto\chi([x,y]).
\]
It is injective in \(x\), with locally split image: the Lie-algebra
factor is the injective principal \(\mathfrak{sl}_2\) operator
\([f_{\mathrm{prin}},-]\) on positive grades, and the coefficient
factor is the perfect residue pairing
\[
 (\overline f,\overline g)\longmapsto
       \operatorname{Res}\frac{f(z)g(z)}{P(z)}\,dz
       \quad\text{on }R[z]/(P).
\]
This pairing is perfect for every monic \(P\), including repeated
roots.  The higher terms in the orbit map at height \(l-1\) involve
only lower-height coordinates.  Induction in height therefore identifies
the orbit with a closed graph over the locally split linear image,
and identifies its parameter space with \(U\).

Closed pushforward from this orbit is fully faithful with continuous
right adjoint, also for categories with coefficients.  On the orbit,
translation descent for the free \(U\)-action identifies equivariant
objects with the fiber over zero.  Under this identification, forgetting
from that fiber and star averaging is exactly the transition at the
break, with the fixed smooth/Fourier normalization.  It is therefore
fully faithful with continuous linear right adjoint.  Composing the
finitely many breaks between successive integral slopes gives the
claimed transition.  After passage to \(b\), such a left adjoint
preserves compact objects because its right adjoint is continuous.  All operations used in the
argument are relative Fourier transform, closed pushforward, smooth
descent, and the base-change operations of Section~\ref{sec:loop};
this proves the remaining assertions.
\end{proof}

\subsection{The point input and its module convention}

We specify the representation-theoretic input before extending it over
\(S\).  For a split root Lie algebra over a characteristic-zero field,
Raskin's affine Skryabin theorem identifies the Whittaker coinvariants
of the renormalized Kac--Moody category with the renormalized category
of principal \(W\)-modules \cite[Theorem~5.1.1]{RaskinW}.
Here the latter category means
\[
 \operatorname{Ind}\bigl(\operatorname{thick}
       \{V_i^W:i\geq0\}\subset D^+(W\text{-}\mathrm{sm})\bigr).
\]
In particular the mapping complexes in its specified compact category
are the ordinary derived mapping complexes between the generalized
vacua.  This is the convention in \cite[Section~5.1 and the proof of
Theorem~5.1.1]{RaskinW}; it does not assert that these objects are
compact in the unrenormalized unbounded derived category.

The composite with the underlying-complex functor is principal
Drinfeld--Sokolov reduction.  It sends the Whittaker averaging
counits to isomorphisms.  On the induced \((I_i,\chi)\)-generator,
using the \(N_i\)-lattice, it gives the generalized vacuum \(V_i^W\)
in degree zero.  Changing the lattice introduces its relative
cohomological shift and determinant line.  The vacuum calculation,
its nonnegative PBW filtration, and the strict surjections
\(V_j^W\twoheadrightarrow V_i^W\) for \(j\geq i\) are
\cite[Theorems~4.2.1, 4.5.1, and 4.6.1]{RaskinW}.
The ordinary vacuum is the case \(i=0\).

We use unshifted vertex modes \(w_{(m)}\).  For homogeneous principal
PBW generators \(w\) of conformal degree \(d_w\), the vacuum is
annihilated by \(w_{(m)}\) for \(m\geq i d_w\), and ordered
monomials in the remaining modes give its PBW basis.  One can recover
the exact, rather than merely associated-graded, annihilation from
the free-field realization \cite[Theorem~6.1.1]{RaskinW}.  The vacuum
embeds in the depth-\(i\) Heisenberg vacuum.  A homogeneous
normally ordered monomial of weight \(d\) has \(k\) currents and
\(r\) derivatives with \(k+r=d\); its mode of index \(m\) has
current indices with sum \(m-(k-1)-r\).  If \(m\geq id\) and
\(i\geq1\), this sum is greater than \(k(i-1)\), so one of the
annihilating current modes occurs on the right after normal ordering.
The creation axiom gives the case \(i=0\).  The associated-graded
calculation then proves the stated PBW basis.

We also use Feigin--Frenkel duality of principal conformal vertex
algebras at the inverse shifted forms.  Our levels are irrational;
thus they lie in the proved range of the duality, including its
vacuum-preserving form \cite[Section~7.3, especially Warning~7.3.1
and Lemma~7.3.3]{RaskinW}.  The input is a Lie-algebra theorem and
is unaffected by the isogeny type of the group.  The principal BRST
construction provides closed fields acting on reduction, with the
usual conformal coordinate action.  One may use its nonnegative
ghost-degree subalgebra: in the principal integer grading it has
cohomology only in degree zero, and its degree-zero kernel realizes
the \(W\)-algebra.  This is the subcomplex \(\mathcal C^-\) of
\cite[Section~4, Theorem~4.1 and Remark~4.1]{KacWakimoto}, with
homological charge reversed to cohomological degree.  Its absence
of negative cohomological degrees means that no degree-zero boundary
has to be quotiented out.  Thus its degree-zero cycles form an actual
vertex subalgebra of the BRST algebra, acting on every module complex.
This supplies the strict action needed here; the corresponding
factorization action is discussed in \cite[Section~4.4]{RaskinW}.

For completeness, the coordinate convention agrees with our
\(P_0\)-twist.  The improved Virasoro term is
\(L_n^{\mathrm{Sug}}-(n+1)\rho_n\), together with its ghost terms.
On a current mode \(x_m\) of height \(l\) its noncentral commutator is
\[
 [L_n,x_m]=-(m+(n+1)l)x_{m+n}.
\]
This is the coordinate action on the root coefficient with density
\((dz)^l\).  The central term is the affine gauge correction for
\((dz)^\rho\); differentiating \(z\mapsto z+\epsilon z^{n+1}\)
gives precisely the correction contributed by \(-(n+1)\rho_n\).
The charged ghosts transform by the dual densities.  Corrected
currents and their normally ordered ghost bilinears preserve the
closed-field model, with possible scalar corrections from reordering.
Consequently the BRST action and the conformal vertex-algebra action
use the same changes of coordinate.

\subsection{Reduction over a moving divisor}

Let \(\mathcal H_R^W\) be the abelian category of smooth modules for
these \(W\)-modes over \(R\).  Concretely, fields are smeared by
coefficients in \(\mathcal K\), with total residue and the translation
and operator-product relations; equivalently one uses the chiral
module construction of \cite[Section~3.6]{BDChiral}.  Smoothness means
that, on any fixed vector, sufficiently high-order coefficients of each generating field
act by zero.  The identity field acts by total residue against \(dz\).
For coefficients \(f,g\), commutators use the vertex products with
coefficients \((\partial^n f/n!)g\).  Composite fields act by their
normal-product formulas.  Write \(D(\mathcal H_R^W)\) for the ordinary
derived category.

These formulas make sense also when the divisor has multiple points.
For example, the two expansions of a pole \((z-w)^{-j}\), \(j>0\),
are obtained by writing
\[
 (z-w)^{-j}
 =\left(\frac{P(z)-P(w)}{z-w}\right)^j(P(z)-P(w))^{-j}
\]
and expanding the last factor in the two orders in the nested
Laurent rings.  In each resulting normally ordered series, the
coefficients of the modes applied first tend to arbitrarily high
order.  Thus each expression acts on a smooth vector by a well-defined
sum.  The construction is unchanged upon replacing \(P\) by an
equation giving a cofinal completion, commutes with scalar extension,
and uses products of mode algebras on disjoint formal divisors.

Give a generating field its principal degree \(d_w\).  A commutator
of two generating modes has strictly smaller total PBW degree: a
nonnegative vertex product has conformal weight smaller than the sum
of the two weights, and expressing it through products and derivatives
does not increase that bound.  We will repeatedly use this fact to
reorder words.  Every such calculation is vectorwise; it requires no
single smoothness bound valid for an entire module.

The reduction functor on compact Kac--Moody modules is the complex
\begin{equation}\label{eq:L6}
 \operatorname{DS}_R(M)=
 C^{\infty/2,h}\bigl(\operatorname{Lie}(LN),M\otimes R_{-\chi}\bigr),
 \qquad h=\operatorname{Lie}N_0(\mathcal O).
\end{equation}
We use cochains on the compact lattice.  If \(h'\supset h\),
changing from \(h'\) to \(h\) tensors the complex by
\(\det(h'/h)[\dim(h'/h)]\).  The critical bracket correction is zero
on the unipotent Lie algebra.  The complex uses the algebraic smooth
module and the Fock ghost module.  On a bounded complex of inputs,
quasi-isomorphisms remain quasi-isomorphisms: the input direction is
bounded and the ghost factors are free.  The compact calculations in
Proposition~\ref{loop:hc-renormalization} therefore define
\eqref{eq:L6} on compact objects, and then define its continuous
extension to \(\mathcal M^0\).

The BRST differential and the closed \(W\)-fields are given by the
same residue formulas on every chart.  Their identities can be
checked in the affine-plus-Clifford enveloping algebra.  The preceding
expansions guarantee convergence in the smooth-module topology;
induced modules for sufficiently deep lattices have free PBW terms,
and the Clifford factor has its free Fock presentation.  These
modules detect the identities.  Equivalently, on them one may check
the identities after every residue-field extension, where they are
the ordinary BRST identities in uniformizers.  This also verifies
closure of the degree-zero fields acting on the complex.

The full regular lattice \(h\), its Fock module, and its vacuum are
canonically identified across nilpotently different lifts of the
parameter scheme.  Indeed the corresponding regular and Laurent
completions are canonically equal.  Thus \eqref{eq:L6}, after forgetting
the \(W\)-action if necessary, defines a crystal functor.  In
particular the ordinary vacuum cycle, consisting of the inducing
vector times the full-lattice Fock vacuum, is horizontal.  We do not
assert that a prescribed positive-depth cutoff is itself horizontally
constant.

\begin{lemma}[Relative depth vacua]\label{walg:vacua}
The reduction of the induced \((I_i,\chi)\)-generator, calculated
with the \(N_i\)-lattice, is a flat heart module \(V_{i,R}^W\).
It has a distinguished vector \(\mathbf v_i\) with presentation
\begin{equation}\label{eq:L7}
 w(u^{i d_w}\mathcal O)\mathbf v_i=0,
 \qquad
 V_{i,R}^W\text{ has PBW basis in modes of orders }<i d_w.
\end{equation}
It corepresents vectors satisfying the displayed annihilations.
These constructions commute with the scalar extensions used here,
including the finite flat infinitesimal tests for the crystal.
The ordinary vacuum satisfies the same assertion for \(i=0\).
\end{lemma}

\begin{proof}
Coefficients of each fixed \(u\)-order have basis
\(z^a\), \(0\leq a<\deg P\).  Apply ordered words in the closed
fields to the inducing vector times the lattice Fock vacuum.
This defines a map from the free \(R\)-module on the proposed PBW
words to the BRST complex.  Its terms and its construction commute
with ordinary scalar extension.  Over every residue field it is a
quasi-isomorphism by the point vacuum calculation: at a point of
multiplicity \(m\), the equation \(u\) has order \(m\), so that
its depth is \(mi\); separated points give the product calculation.

The BRST terms are flat over \(R\).  Since \(R\) is smooth of finite
dimension, the scalar extensions to residue fields have finite Tor
dimension, and ordinary tensor with these complexes computes derived
tensor.  The cone of the displayed map therefore has zero derived
fiber at every prime.  Residue fields at all primes detect zero
objects of \(D(R)\), so the map is a quasi-isomorphism.  Its
cohomology is the indicated free module.  The annihilations follow
from the point annihilations, since an operator between free
\(R\)-modules that is zero at every prime is zero.  The same free
presentation remains valid on the finite flat infinitesimal tests.

Finally let \(v\) be a smooth vector satisfying \eqref{eq:L7}.
Reorder a word by decreasing total PBW degree, moving annihilating
modes to the right; each commutator lowers degree.  The same operation
on \(\mathbf v_i\) gives its PBW expansion.  At each comparison the
normal-product tails terminate by the smoothness bounds for the two
vectors.  Hence the substitution \(\mathbf v_i\mapsto v\) respects
every mode action.  Its uniqueness follows from cyclicity, proving
the universal property.
\end{proof}

The modules \(V_{i,R}^W\) generate \(\mathcal H_R^W\), because every
vector satisfies a sufficiently deep cutoff.  Kernels, cokernels,
and filtered colimits are computed on the underlying modules, with
the induced smooth action.  Consequently \(\mathcal H_R^W\) is a
Grothendieck category with exact filtered colimits.

\subsection{Mapping complexes and scalar extension}

The passage from the point theorem to families needs a finiteness
statement for derived maps.  Smoothness alone would not justify
commuting an infinite collection of mode conditions with tensor
product.  The finite differences between the depth vacua provide
the required replacement.

\begin{lemma}[Ext continuity]\label{walg:ext}
For every depth \(i\geq0\) and every \(p\geq0\), the functor
\(\operatorname{Ext}^p_{\mathcal H_R^W}(V_{i,R}^W,-)\) commutes
with filtered colimits.  If \(R\to T\) is a scalar extension occurring
in our affine or infinitesimal charts, then, for a smooth \(W\)-module
\(M\) over \(T\), restriction gives
\begin{equation}\label{walg:restriction-ext}
 \operatorname{Ext}^p_{\mathcal H_T^W}(V_{i,T}^W,M)
 \simeq
 \operatorname{Ext}^p_{\mathcal H_R^W}
             (V_{i,R}^W,\operatorname{oblv}_R M).
\end{equation}
For smooth affine or residue-field extensions, and for the finite
flat infinitesimal tests, the ordinary derived mapping complexes
between depth vacua consequently satisfy
\begin{equation}\label{walg:mapping-base-change}
 \operatorname{RHom}_{\mathcal H_R^W}(V_{i,R}^W,V_{j,R}^W)
       \otimes_R^{\mathbf L} T
 \simeq
 \operatorname{RHom}_{\mathcal H_T^W}(V_{i,T}^W,V_{j,T}^W).
\end{equation}
\end{lemma}

\begin{proof}
\emph{Resolving the difference between two depths.}
Put \(K_{ji}=\ker(V_{j,R}^W\twoheadrightarrow V_{i,R}^W)\) for
\(j>i\).  We first construct a resolution of \(K_{ji}\), possibly
infinite to the left, whose terms are finite sums of depth vacua.
The PBW filtration on its defining surjection is strict; in its
associated graded the kernel is generated by the finitely many mode
coefficients between the two cutoffs.  Choose homogeneous generators
and lift them to \(K_{ji}\).  Their smoothness supplies a common
sufficiently deep vacuum mapping to each chosen lift, with the
appropriate filtration shift.  These maps give a strict surjection
onto \(K_{ji}\).

The construction can be iterated.  Indeed, at each step choose a
common upper depth \(m\).  All the associated-graded modules involved
are finitely presented over the polynomial ring in the mode variables
below depth \(m\).  Modes above that depth act as zero in associated
graded, because commutators lower degree; they need not annihilate the
filtered modules.  Increasing from one finite depth to another adds
only finitely many variables.  The polynomial ring in all the
remaining, countably many, variables over the Noetherian ring \(R\)
is coherent.  To verify precisely the fact used here, a matrix between
finite free modules and its finitely many coefficients descend to a
polynomial ring in finitely many variables; its kernel is finitely
generated there, and extension to the full polynomial ring is flat.
Thus kernels between finitely presented modules remain finitely
presented.  The next graded kernel is therefore finitely generated,
and the lifting argument repeats.  This proves the resolution
assertion.

Every vacuum is \(R\)-flat.  The first kernel \(K_{ji}\) is flat
because it is the kernel of a surjection between flat modules with
flat quotient.  Inductively, every subsequent kernel in the resolution
is flat for the same reason.  Hence these exact sequences remain
exact after arbitrary scalar extension in the charts.  The PBW
presentations of Lemma~\ref{walg:vacua} identify the extended sequences
with the corresponding sequences over \(T\).  This argument applies
also over the finite flat infinitesimal charts.

\smallskip
\noindent\emph{Continuity of Ext under filtered colimits.}
For \(p=0\), a map from a vacuum is determined by its generating
vector.  If that vector is given in a filtered colimit, lift it to
one member.  The lift is annihilated above some depth \(j\).
Only finitely many coefficients between depth \(i\) and depth \(j\)
remain to be tested.  Their vanishing holds in a common later member,
which proves continuity of \(\operatorname{Hom}(V_{i,R}^W,-)\).
The universal property also gives \eqref{walg:restriction-ext} for
\(p=0\).

For the induction in \(p\), the exact underlying-module functor has
the presentation
\[
 \operatorname{oblv}(M)=
 \underset{j}{\operatorname{colim}}
       \operatorname{Hom}(V_{j,R}^W,M).
\]
Applying this formula to an injective resolution and using exactness
of filtered colimits gives
\[
 \underset{j}{\operatorname{colim}}
       \operatorname{Ext}^p(V_{j,R}^W,M)=0\qquad(p>0).
\]
The long exact sequences for \(K_{ji}\), followed by this colimit,
therefore give, for every \(p>0\), the useful formula
\begin{equation}\label{walg:ext-bootstrap}
 \operatorname{Ext}^p(V_{i,R}^W,M)
 \simeq\operatorname{coker}\left(
 \underset{j>i}{\operatorname{colim}}
       \operatorname{Ext}^{p-1}(V_{j,R}^W,M)
 \longrightarrow
 \underset{j>i}{\operatorname{colim}}
       \operatorname{Ext}^{p-1}(K_{ji},M)\right).
\end{equation}
For fixed \(j,i\), the resolution just constructed calculates
\(\operatorname{Ext}^{p-1}(K_{ji},M)\) by the first-quadrant
resolution spectral sequence.  Its terms are finite sums of
\(\operatorname{Ext}^q(V_{m,R}^W,M)\), and its diagonals in any
fixed total degree are finite.  Only \(q\leq p-1\) is needed.
The induction hypothesis therefore gives filtered-colimit continuity
for this Ext group.  Formula~\eqref{walg:ext-bootstrap} proves the
claim in degree \(p\).

\smallskip
\noindent\emph{Restriction under scalar extension.}
The same induction proves \eqref{walg:restriction-ext}.  Its degree
zero assertion was proved above; use the base-extended resolutions
of \(K_{ji}\), which remain exact, in the spectral sequences on the
two bases, and then use \eqref{walg:ext-bootstrap}.  All comparisons
are the natural maps obtained by restriction and the canonical
\(V_{i,R}^W\to V_{i,T}^W\), so this is an identification of the
natural derived mapping complexes, not just a coincidence of their
cohomology dimensions.

\smallskip
\noindent\emph{Derived tensor base change.}
It remains to justify tensoring the target in
\eqref{walg:mapping-base-change}.  First tensor by a finite projective
\(R\)-module; the assertion follows from finite additivity and
retracts.  A flat module is a filtered colimit of finite free modules,
so the Ext-continuity statement proves the assertion for a flat
module as well.  The extensions from a smooth affine \(R\) under
consideration have finite Tor dimension; resolve \(T\) by a bounded
complex of flat \(R\)-modules.  The preceding comparisons totalize
over this finite range and prove the required derived tensor
identity.  Formula~\eqref{walg:restriction-ext} then changes the
ambient module category from \(R\) to \(T\).  Finite flat
infinitesimal tests are already covered by the flat case.
\end{proof}

\subsection{The relative Skryabin comparison}

The stages in \eqref{eq:L5} have positive depth, so we first use
only their vacuum images.  Let \(\mathcal V_R^W\) denote the thick,
idempotent-complete subcategory of \(D(\mathcal H_R^W)\) generated
by \(V_{i,R}^W\), \(i\geq1\).  Keeping the ordinary mapping
complexes in this definition is essential.  After constructing the
relative comparison, Lemma~\ref{walg:ordinary-vacuum} will prove
geometrically that adjoining \(V_{0,R}^W\) does not change this
category.  Define the relative renormalized category by
\[
 \mathcal W_R^{\mathrm{ren}}=
       \operatorname{Ind}(\mathcal V_R^W).
\]
The next proposition constructs this category intrinsically as the
Whittaker target and then checks its dependence on the formal curve.

\begin{proposition}[Relative Skryabin comparison]\label{walg:relative-skryabin}
Reduction induces an equivalence
\[
 Q_S^0\simeq\mathcal W_R^{\mathrm{ren}}.
\]
On compact objects its composite with the ordinary comparison to
\(D(\mathcal H_R^W)\) is \eqref{eq:L6}.  This comparison is fully
faithful on compact objects and has essential image
\(\mathcal V_R^W\).  It commutes with the scalar extensions in
Lemma~\ref{walg:ext} and with the crystal transport of
Section~\ref{sec:loop}.
\end{proposition}

\begin{proof}
We first descend the continuous extension of \eqref{eq:L6} through
the Whittaker projection.  Apply it to the counit of star averaging
by a compact subgroup \(N_i\).  After forgetting the \(W\)-action,
this comparison commutes with tensor to every residue field:
on congruence compact generators this follows from the BRST formula,
and on their continuous extensions it follows by colimits.  The
point Skryabin theorem says that every such fiber of the counit is
an isomorphism.  Residue-field detection in \(D(R)\) therefore makes
the original counit an isomorphism.  The isomorphisms are the images
of the averaging counits themselves and hence carry their composition
coherences.  They define a single functor
\[
 Q_S^0\longrightarrow D(\mathcal H_R^W)
\]
using \eqref{eq:L5}, whose composite with the canonical projection
from \(\mathcal M^0\) is reduction.  In particular this is a
comparison on functors and maps, as well as on objects.

By Proposition~\ref{walg:adolescent}, every compact object of
\(Q_S^0\) is a retract of an object generated at a finite stage,
and the stage transitions preserve compactness.  Each stage is
thickly generated over \(R\) by the induced trivial representation,
with its character correction.  This is the pro-unipotent
Harish--Chandra calculation of
Proposition~\ref{loop:hc-renormalization}.  Its reduction is the
vacuum of Lemma~\ref{walg:vacua}, with the prescribed relative lattice
line and shift.

Take two such generators.  Move them to a common later stage of
\eqref{eq:L5}.  Their source mapping complex commutes with the allowed
scalar extensions by Proposition~\ref{loop:hc-renormalization}
and the continuous linear averaging adjunctions.  Their target
mapping complex commutes with these extensions by
Lemma~\ref{walg:ext}.  On each residue field the resulting mapping
comparison is the point Skryabin equivalence.  Apply residue-field
detection to its cone.  The comparison is consequently an
isomorphism over \(R\), proving full faithfulness on the generators
and hence on their thick closure.  Its essential image is exactly
\(\mathcal V_R^W\); taking Ind proves the equivalence.

For crystal transport, work on a finite infinitesimal neighborhood
of a diagonal in the parameter scheme.  The two lifts have the same
absolute completed curve, full regular lattice, and Laurent algebra.
The ordinary BRST comparison on Kac--Moody compact objects is
therefore the same for both lifts, including its Fock vacuum.
Its mapping comparisons agree by the finite flat case of
Lemma~\ref{walg:ext}.  To pass to \(Q\), refine to common deep
congruence subgroups as in Proposition~\ref{loop:crystal}; their
systems are cofinal.  On an equivariant stage the projection to
coinvariants is precisely its map into the colimit, so its ordinary
comparison is calculated by first forgetting to the Kac--Moody
category and applying the same BRST functor.  Thus the comparisons
agree on the stages and their transition maps, and hence on the
colimit.  The identifications are induced by the identity of the
absolute formal curve.  They compose identically on higher
infinitesimal diagonals, which supplies the crystal coherence.
\end{proof}

\begin{lemma}[The ordinary vacuum is in the compact image]
\label{walg:ordinary-vacuum}
The image in \(Q_S^0\) of the spherical Kac--Moody vacuum is compact.
Its ordinary comparison is \(V_{0,R}^W\).  Consequently adjoining
\(V_{0,R}^W\) to the generators of \(\mathcal V_R^W\) does not
change that category.
\end{lemma}

\begin{proof}
Let \(\mathrm{Vac}_H\in C_S\) be the module induced from the
trivial arc-group representation, and let \(\delta_1\in a_S\)
be the convolution unit supported at the unit section of the
Grassmannian.  By \eqref{eq:L2}, \(b\mathrm{Vac}_H\) is compact
in \(C_S^0\).  We first verify that
\(b\mathbf d(\delta_1)\) is compact in \(d_S^0\).

In the global pole model of Section~\ref{sec:whittaker}, the zero-pole
bound has just one relevant label, namely zero.  The stratum tests
remove all defective strata, leaving the fixed stack
\(\operatorname{Bun}_{N_0}\times S\).  Exponential pullback identifies
its Whittaker category with \(D(S)\), with the fixed determinant line
and normalization from that section.  This assertion is uniform in
the marks: a zero divisor does not change when marks collide.
The inclusion of this Whittaker pole bound into the full pole model
has continuous \(D(S)\)-linear right adjoint given by \(!\)-restriction,
by Propositions~\ref{whittaker:strata-properties} and
\ref{whittaker:duality}.  Tensoring this adjunction with
\(\operatorname{QCoh}(S)\) shows that the image of the compact
module \(R\) is compact after passage to \(b\).

Under \eqref{eq:W2}, this opposite zero-label object represents
\(\mathbf d(\delta_1)\), up to the fixed invertible line and shift
specified by the pairing.  To verify the identification, its
functional on a Whittaker object \(W\) is the normalized zero-label
coefficient.  Indeed restriction to the zero-pole bound extracts that
coefficient by exponential cancellation.  The functional defining
\(\mathbf d(\delta_1)\) in Section~\ref{sec:loop} is
\[
 W\longmapsto
 \Gamma_{\mathrm{dR}}(\delta_1\otimes^!\mu^!W).
\]
It is the same coefficient: the Grassmannian unit maps to the trivial
\(N_0\)-torsor, a section of the zero-label stratum over \(S\),
and normalized \(!\)-pullback to that section inverts exponential
pullback of the coefficient.  More explicitly, if \(q_0\) is the
zero-label projection and \(\sigma\) the trivial-torsor section,
then \(q_0\sigma=\operatorname{id}_S\) and
\(\psi_0\sigma=0\).  On the pairing side,
\(q_{0,*}q_0^!\) contributes \([2\delta_N]\), exactly canceled
by the \([-2\delta_N]\) in \eqref{eq:W2}.
Proposition~\ref{whittaker:duality}
identifies closed inclusion with the transpose of \(!\)-restriction,
so this identifies the two functionals on every object, not just
on standard generators.  Invertible lines and shifts preserve
compactness.  This proves the assertion about
\(\mathbf d(\delta_1)\).

By Proposition~\ref{loop:pairing-adjoints}, the pairing \(F_S\) has a
continuous \(D(S)\)-linear right adjoint, as the composite of the
adjoints to balancing and to the balanced pairing.  After passage to
\(b\) its adjunction remains continuous and linear.  Thus it
preserves compact objects.  The exterior product of the two compact
inputs above is compact, and the convolution unit gives
\[
 F_S(\mathrm{Vac}_H,\mathbf d(\delta_1))
      =q(\mathrm{Vac}_H),
\]
where \(q:\mathcal M\to Q_S\) is the coinvariant projection.
Hence this particular coinvariant image is compact; no compactness
claim about \(q\) on arbitrary inputs is being used.

Proposition~\ref{walg:relative-skryabin} identifies its ordinary
comparison with \(\operatorname{DS}_R(\mathrm{Vac}_H)\), which is
\(V_{0,R}^W\) by Lemma~\ref{walg:vacua}.  The already proved
full faithfulness on compact objects has essential image
\(\mathcal V_R^W\), generated by positive depths.  Therefore
\(V_{0,R}^W\) belongs to that image, as required.
\end{proof}

The relative renormalization thus agrees with the convention that
includes the ordinary vacuum.  The order of the argument matters:
its compactness follows from the geometric pairing, not from the
possibly infinite resolution in Lemma~\ref{walg:ext}.

\subsection{Duality and compatibility}

For either Lie algebra, compose the ordinary comparison of
Proposition~\ref{walg:relative-skryabin} with the underlying-complex
functor:
\[
 Q_S^0\longrightarrow D(\mathcal H_R^W)\longrightarrow D(R).
\]
Its crystal compatibility descends this composite to a continuous
\(D(S)\)-linear functor
\begin{equation}\label{eq:walg-forgetful}
 \operatorname{oblv}_W:Q_S\longrightarrow D(S).
\end{equation}
After passage to \(b\), its composite with the coinvariant projection
is the continuous Drinfeld--Sokolov functor computed on compact
inputs by \eqref{eq:L6}.  The scalar-extension comparisons used in
this construction are those of Lemma~\ref{walg:ext}: finite-Tor
smooth-affine and residue-field extensions, and finite flat
infinitesimal tests for the crystal.

\begin{theorem}[Common Whittaker target]\label{walg:common-target}
For the inverse shifted levels on the two Langlands-dual Lie
algebras, there is a canonical identification of the crystal
categories \(Q_S\).  It respects pullback in \(S\), collisions of
marks, changes of curve coordinate, and exterior products for
disjoint clusters.  The underlying-complex functors from these
categories to \(D(S)\) are identified.  The ordinary vacuum and its
distinguished horizontal generating cycle are identified as well.
\end{theorem}

\begin{proof}
The conformal Feigin--Frenkel isomorphism identifies the two
\(W\)-algebras \cite{FeiginFrenkel}.  At a point it identifies depth
vacua with their cyclic vectors, by \cite[Lemma~7.3.3]{RaskinW}.  In the relative
presentation \eqref{eq:L7} this identification is equally explicit.
The principal degrees agree.  Replacing homogeneous principal PBW
generators by the images of the dual generators does not change the
presentation.  We verify the needed cutoff assertion before applying
its universal property.  Suppose \(i\geq1\), and a vector \(v\) obeys
the generating-field cutoffs.  Induction on conformal weight proves
that every homogeneous field of weight \(d\) has modes of index
\(m\geq id\) annihilating \(v\).  Derivatives preserve the assertion:
\((\partial^r A)_{(m)}\) is a scalar multiple of \(A_{(m-r)}\),
and \(m\geq i(a+r)\) implies \(m-r\geq ia\) for a field of
weight \(a\).

For a normal product of fields of positive weights \(a,b\), use
\[
 (:AB:)_{(m)}=
 \sum_{k<0} A_{(k)}B_{(m-1-k)}+
 \sum_{k\geq0}B_{(m-1-k)}A_{(k)}.
\]
If \(m\geq i(a+b)\), the first sum kills \(v\), as do the terms
\(k\geq ia\) in the second.  In the remaining finitely many terms,
commute \(B_{(s)}\) past \(A_{(k)}\), where \(s=m-1-k\geq ib\).
The reordered term kills \(v\).  Each commutator term is a scalar
multiple of \((B_{(n)}A)_{(m-1-n)}\); this field has weight
\(c=a+b-n-1<a+b\), and
\[
 m-1-n-ic=(m-i(a+b))+(i-1)(n+1)\geq0.
\]
It therefore kills \(v\) by induction.  Weight-zero fields contribute
only the identity mode of index \(-1\), which cannot occur in this
range.  Strong generation now proves the assertion.  The case
\(i=0\) is the ordinary creation axiom.

Over the moving divisor, specialize the transported generator action
to every residue field.  The calculation just made applies at each
support point, with its multiplicity multiplying the depth.  Thus
these modes annihilate the cyclic vector on every fiber.  The module
\(V_{i,R}^W\) is \(R\)-free and \(R\) is reduced, so the same
annihilation holds over \(R\); it persists under the flat
infinitesimal base changes by its presentation.  Hence the images of
the dual generators satisfy \eqref{eq:L7}.  Applying the inverse change of
generators proves equality of the two annihilation conditions.
The common cyclic vector and Lemma~\ref{walg:vacua} determine mutually
inverse maps of the relative vacua.

It follows that Feigin--Frenkel identifies the two thick subcategories
\(\mathcal V_R^W\), with their ordinary mapping complexes, and hence
their Ind completions.  Proposition~\ref{walg:relative-skryabin}
identifies both with \(Q_S^0\).  Its crystal compatibility identifies
this equivalence across the formal diagonals, so it descends to the
claimed equivalence of crystal categories \(Q_S\).

The coordinate compatibility follows from the conformal coordinate
action on the strict closed fields and the BRST calculation above.
For disjoint divisors the modes, Fock modules, and semi-infinite
complexes are tensor products.  The determinant lines and shifts use
the graded tensor symmetry, so the comparisons also respect products
and their permutation coherences.  Finally, the Feigin--Frenkel
identification changes the algebra acting on a module and preserves
its underlying complex.  The two ordinary comparison functors to
\(D(S)\) consequently agree.  At depth zero
Lemma~\ref{walg:ordinary-vacuum} places both vacuum images in the compact subcategories where the ordinary comparison is
fully faithful.  The vacuum-preserving module isomorphism therefore
lifts to their unique isomorphism in \(Q_S^0\), and its crystal
compatibility follows from the same comparison on infinitesimal
charts.  It preserves the actual inducing/Fock cycle, which was
shown to be horizontal.
\end{proof}

Theorem~\ref{walg:common-target} supplies the common target for the
pairings of \eqref{eq:L3}, with the relative structure needed to
compare their right adjoints.  Its proof has used the point
representation theorems together with explicit relative mode and
mapping calculations; no factorization form of a fundamental local
equivalence enters this construction.
